\chapter{Causal Graph and Effect-Estimation Methodology}
\label{chap:causal-graph-effect-estimation-methodology}

\section{Dataset-Specific DAGs}
\label{sec:causal-graph-effect-estimation-methodology:dataset-specific-dags}
The causal component of this project does not learn a causal graph from the observational data.  It uses two LLM-assisted, project-specified directed acyclic graphs (DAGs), one for PhysioNet 2012 and one for MIMIC-III, as explicit design assumptions for selecting adjustment variables and organizing later effect-estimation outputs.  A structured clinical and causal elicitation process generated graph proposals that were selected and encoded as active graph-construction source code.  Work on LLM-derived soft priors provides a methodological precedent for using semantic judgments as candidate causal relations or orientation information \cite{darvariu_et_al_2024_llm_causal_graph_priors}, but it does not validate these local graphs or imply that its discovery algorithm was implemented here.  The diagrams should therefore be read as source-encoded causal hypotheses, not as clinically validated or statistically discovered ground truth.

This separation is important because the main exposure variables in the causal stage are binary proxy states rather than randomized treatments.  In the causal code, each selected proxy state is treated as a binary exposure \(A\), the outcome is in-hospital mortality \(Y\), and the graph is used to identify candidate pre-exposure background or latent proxy variables for adjustment.  The thesis follows the graphical language of Pearl's DAG and backdoor framework while retaining the cautions around observational, proxy-defined exposures and well-defined interventions discussed by Hernan and Taubman \cite{pearl_1995_causal_diagrams,hernan_taubman_2008_well_defined_interventions}.  These cautions are especially relevant here because several variables named as ``treatments'' in the code are clinically closer to patient states or care-process proxies than to assigned therapeutic interventions.

The LLM has no estimator-time role in this chapter's workflow.  It is not called to inspect patient records, orient edges dynamically, select adjustment variables, or choose estimators during execution.  The active graph scripts define the implemented nodes and edges, and the matching and CATE source code determines the implemented adjustment workflow from those encoded graphs and the columns available in each analysis dataframe.

\begin{figure}[p]
\centering
\includegraphics[width=0.95\textwidth]{physionet_causal_dag.png}
\caption{Project-specified PhysioNet 2012 DAG used by the causal pipeline.  The active graph source reproduces the archived artifact; arrows encode assumed rather than learned or clinically validated relationships.}
\label{fig:physionet-causal-dag}
\end{figure}

\begin{figure}[p]
\centering
\includegraphics[width=0.95\textwidth]{mimic_causal_dag.png}
\caption{Project-specified MIMIC-III DAG used by the causal pipeline.  The active graph source reproduces the archived artifact; arrows encode assumed rather than learned or clinically validated relationships.}
\label{fig:mimic-causal-dag}
\end{figure}

Both graphs distinguish background variables, proxy states, observed measurements, care processes, missingness or measurement-intensity processes, and the mortality outcome.  Only mapped background and eligible upstream proxy nodes can enter adjustment sets; measurement, care-process, missingness, and outcome nodes remain in the causal hypothesis but are excluded as adjustment candidates.  MIMIC includes richer administrative and care-process structure, whereas PhysioNet uses a smaller background and observation vocabulary.

The PhysioNet graph contains 30 nodes and 45 directed edges.  The implemented exposure list contains the ten non-chronic latent proxy states in the PhysioNet configuration: global severity, shock, respiratory failure, renal dysfunction, hepatic dysfunction, coagulation/hematologic dysfunction, inflammation/sepsis burden, neurologic dysfunction, cardiac injury/strain, and metabolic derangement.  Chronic baseline risk is retained as an upstream baseline proxy and effect-modifier candidate rather than as an analyzed exposure.

The MIMIC graph contains 36 nodes and 57 directed edges.  The implemented exposure list contains nine non-chronic latent proxy states: global severity, inflammation/sepsis, shock, respiratory failure, renal dysfunction, hepatic/coagulation dysfunction, neurologic dysfunction, cardiac strain, and metabolic derangement.  Chronic burden is retained as an upstream baseline proxy and effect-modifier candidate.  The MIMIC graph contains richer administrative and care-process structure than the PhysioNet graph, but the downstream adjustment implementation only uses graph nodes that can be mapped into the analysis dataframe.

Qualified clinical review of the graph arrows and a complete record of the human decisions applied to the LLM-assisted proposals remain unavailable.  The graph artifacts and active source structure are recorded, but the exact producing source version, configuration, and command are not fully recovered.

The archived final run summaries and majority-vote logs establish the downstream proxy-label lineage at the stage boundary.  In all twelve runs, the vote combines five enumerated inputs: one deterministic rule-derived proxy-label table and four thresholded prediction tables from GRU, GRU-D, STraTS, and TCN.  Each run writes \nolinkurl{majority_vote/latent_tags_majority_vote.csv} and passes that mixed rule-and-prediction aggregate to mortality prediction, matching, and CATE estimation.  The logs preserve the external absolute filenames, but not the voter-file bytes, hashes, or checkpoint-to-export mappings; those parts of the lineage therefore remain incomplete.

\section{Adjustment-Set Logic}
\label{sec:causal-graph-effect-estimation-methodology:adjustment-set-logic}
For each dataset, treatment-like proxy state, and outcome pair, the causal pipeline derives an adjustment set from the project DAG and the columns available in the analysis dataframe.  The implementation first maps dataframe variables to graph nodes and graph nodes back to dataframe columns.  Background and latent graph nodes are eligible only if they can be represented by available columns.  Observed measurement nodes, care-process nodes, missingness nodes, and the outcome node are excluded from candidate adjustment sets.

Let \(G\) denote the implemented project DAG, \(A\) the binary proxy-state exposure, and \(Y\) the in-hospital mortality outcome.  The code constructs a graph \(G_{\underline{A}}\) by removing outgoing edges from \(A\).  This is a project helper for finding adjustment candidates; it should not be confused with the usual intervention graph notation that removes incoming edges into \(A\).  The initial candidate pool is:

\[
\mathcal{C}(A,Y)=
\mathcal{A}_{\mathrm{allowed}}
\cap \operatorname{An}_{G}(A)
\cap \operatorname{An}_{G_{\underline{A}}}(Y)
\setminus \operatorname{De}_{G}(A)
\setminus \{A,Y\},
\]

where \(\mathcal{A}_{\mathrm{allowed}}\) is the set of available graph nodes whose node type is background or latent.  The implementation then enumerates backdoor paths between \(A\) and \(Y\) in the undirected skeleton whose first directed edge points into \(A\).  A path is treated as blocked if an adjusted non-collider is present or if an unadjusted collider has no adjusted descendant.  This is an encoded graphical test over the project DAG, not proof that all real-world confounding has been controlled.

The preferred helper is the safe expanded adjustment mode.  It removes encoded colliders and their descendants from the full candidate pool, then tests whether the remaining mapped variables block every encoded backdoor path.  If not, a deterministic greedy fallback removes sorted candidates only while path blocking is preserved.  The returned set is inclusion-minimal under that order, not necessarily unique.  The accompanying identifiability flag is therefore an availability diagnostic within the implemented graph, not proof of exchangeability in the clinical population.

\begin{longtable}{@{}L{0.16\textwidth}L{0.25\textwidth}L{0.22\textwidth}L{0.27\textwidth}@{}}
\caption{Causal assumptions required for interpreting the estimator outputs.}
\label{tab:causal-assumptions}\\
\toprule
assumption & methodological requirement & repository support & unresolved limitation \\
\midrule
\endfirsthead
\toprule
assumption & methodological requirement & repository support & unresolved limitation \\
\midrule
\endhead
Graph correctness for the target question & The DAG must correctly represent the relevant causal structure for the exposure contrast and mortality outcome. & Dataset-specific graphs were developed through LLM-assisted elicitation, project selection, and active source encoding. & The arrows were not learned from data, and the human clinical review record remains incomplete. \\
Sufficient measured adjustment & Available adjustment variables must block the relevant open backdoor paths for each proxy exposure. & The helper tests path blocking within the implemented DAG using mapped background and latent columns. & Graph variables may be unavailable, and unmeasured or misspecified confounding outside the project DAG remains possible. \\
Positivity or overlap & Comparable covariate strata must contain both exposed and unexposed records. & Matching outputs report match availability and match rates. & Dedicated treatment-specific overlap diagnostics and thresholds remain to be approved. \\
Consistency and well-defined exposure & Each binary proxy-state contrast must have a sufficiently stable meaning across records. & The pipeline applies one binary exposure column at a time. & Disease-state and severity proxies are not necessarily assignable interventions, so the causal estimand remains unresolved. \\
Proxy-state measurement & The binary exposure should measure the intended construct with sufficiently limited construct and classification error. & Final archived analyses use a deterministic aggregate of one rule-derived source and four predicted-label sources trained against the rule-derived targets. & No code path verifies construct validity or removes proxy-state measurement error; clinical and chart-review validation remains incomplete. \\
Outcome measurement & The processed \nolinkurl{in_hospital_mortality} field must be consistently and correctly constructed from the outcome artifact. & Both estimator paths merge this processed outcome column by record identifier and require it to be present. & Source checks establish schema availability, not endpoint validity or cross-dataset construction equivalence. \\
Temporal ordering and avoidance of post-exposure adjustment & Adjustment variables must be measured or conceptually occur before the relevant proxy-state contrast; post-exposure adjustment should be avoided. & Protection depends on the project DAG, source node classification, descendant filtering, and allowed-node filtering. & Repository graph ordering does not prove patient-record timing; actual variable timing and the correctness of the proxy-state construction window remain unverified. \\
Nuisance- and effect-model adequacy & LinearDML and CausalForestDML require sufficiently adequate nuisance and effect models, regularity conditions, and appropriate cross-fitting behavior for their intended interpretation. & The source configures EconML estimators, random-forest nuisance models, and estimator-provided fitting behavior. & Source configuration alone does not verify model adequacy, regularity conditions, calibration, or valid cross-fitting performance in these data. \\
No interference & One record's proxy-state exposure should not affect another record's outcome under the working unit-level interpretation. & Estimation treats records as separate units. & The repository does not empirically test interference or dependence across admissions, patients, units, or care processes. \\
Sampling and target population & The analysis sample must support the stated target population, or reweighting or another transport argument must justify a sampled target. & The optional outcome-downsampling condition is deterministic given the configured seed and is applied consistently before treatment iteration. & Original and outcome-downsampled analyses correspond to different empirical populations; no weighting or target-population correction is documented. \\
\bottomrule
\end{longtable}

\section{Matching and Heterogeneous-Effect Estimators}
\label{sec:causal-graph-effect-estimation-methodology:matching-heterogeneous-effect-estimators}
The causal pipeline applies two families of estimators to the mixed-source majority-vote proxy-label exposures: a transparent matched-pair baseline and model-based heterogeneous-effect estimators.  Both families use the same broad analysis frame.  Patient-level aggregated proxy labels are merged with the canonical in-hospital mortality outcome and baseline covariates.  For a selected exposure \(A_i \in \{0,1\}\), outcome \(Y_i\), adjustment variables \(W_i\), and effect modifiers \(X_i\), the target working quantity is a contrast in predicted or matched mortality outcome between \(A_i=1\) and \(A_i=0\), conditional on the source-defined covariate information.  Because the exposures are proxy states and the data are observational, these contrasts require cautious interpretation.

Measurement error can propagate through this entire chain.  An unsupported design proposal can enter a selected rule or edge; a deterministic rule can misclassify the intended construct; a deep model can reproduce that error or add prediction error; aggregation can preserve correlated error; and the resulting exposure misclassification can alter adjustment, matching, and effect estimates.  An LLM-assisted graph is not intrinsically more correct than a manually proposed graph, and predictive agreement with source-encoded labels does not repair graph or exposure-measurement assumptions.

An optional analysis condition, controlled by \texttt{DOWN\_SAMPLE}, down-samples mortality outcome classes rather than exposure groups.  All outcome-positive rows are retained.  When the outcome-negative class is larger and the positive class is nonempty, outcome-negative rows are sampled without replacement until their count equals the number of outcome-positive rows; the configured random seed controls both this sample and the subsequent row shuffle.  This occurs once before treatment-specific adjustment selection and estimation, including matching and DML/PFN fitting.  It is not propensity trimming, treatment balancing, matching, or a correction for confounding, and it does not improve causal identification by itself.  The procedure changes the analyzed population and observed outcome prevalence.  The thesis therefore treats original cohorts as primary and outcome-downsampled cohorts as robustness analyses only, with no pooling between them.

Outcome-based downsampling can change nuisance-model estimation, the distribution of covariates and proxy exposures, patient-level fitted effects, the population over which \texttt{mean\_cate} is averaged, and the availability and composition of matched pairs.  It also changes the sample outcome rate and therefore affects \texttt{normalized\_CATE}, whose denominator is that rate.  Effect summaries from original and outcome-downsampled analyses should not be compared as though they estimated the same population quantity without an approved sampling and estimand argument.  Outcome-based case-control sampling does not necessarily invalidate every estimator, but interpretation here requires estimator-specific methodological justification and an explicit target-population argument.

The matching estimator is a deterministic baseline.  For each proxy exposure, the pipeline drops rows with missing exposure or outcome, coerces the exposure and outcome to binary variables, imputes numeric adjustment variables with medians, and converts non-binary numeric adjustment variables into binary columns using sample-median thresholds.  Already-binary variables, including ICU-type indicators, remain binary.  The matcher then greedily pairs exposed and unexposed patients by Hamming distance over these binary adjustment columns, increasing the allowed distance up to the configured maximum.  By default, controls are not reused.  The per-pair contrast is

\[
\Delta_j = Y_{j,A=1} - Y_{j,A=0},
\]

and the matching summary reports the mean and standard deviation of these matched-pair differences, the number of matched pairs, the match rate, the final allowed Hamming distance, and the observed adjustment columns.  This output is useful as a transparent matched descriptive contrast.  It should not be labeled as an average treatment effect unless the estimand, graph assumptions, overlap, and matching design are later approved for that interpretation.

For model-based heterogeneous effects, the pipeline includes EconML double-machine-learning estimators and an exploratory CausalPFN branch.  The prespecified hierarchy treats CausalForestDML as primary, LinearDML as the secondary comparator, and CausalPFN as exploratory.  Double machine learning estimates nuisance outcome and treatment models before estimating effects from residuals, reducing sensitivity to nuisance-model overfitting under appropriate assumptions \cite{chernozhukov2018dml,oprescu_et_al_2019_econml}.  The LinearDML branch uses random-forest nuisance models with a linear final effect model.  The CausalForestDML branch uses the same nuisance-model family and a causal-forest final stage, drawing on generalized random-forest and causal-forest ideas for heterogeneous treatment effects \cite{wager2018causalforest,athey2019grf}.  This choice is consistent with broader machine-learning work on individualized treatment effects in electronic health records and critical care, while still requiring domain-specific validation \cite{smit_2023_causal_inference_icu_scoping_review,bica_2021_individualized_treatment_effects_ehr_ml,lipkovich_2024_modern_hte_methods,curth_2024_ml_individualized_treatment_effects,iwashyna_2015_hte_critical_care}.

The CATE pipeline saves patient-level CATE values, normalized CATE values, optional 95 percent effect-interval columns when the estimator provides them, feature-importance or coefficient diagnostics when available, model artifacts, and run-level summaries.  In the current source, \(\texttt{mean\_cate}\) is the arithmetic mean of patient-level CATE estimates over the model rows.  It is an aggregate of the fitted model's conditional estimates and should not automatically be described as the study ATE.  The normalized CATE column is computed by dividing CATE by the sample outcome rate when that rate is positive.  It is an implementation-defined rescaling for comparison within this project; it is not a risk ratio, relative risk, or percent reduction.

The effect-modifier matrix \(X\) is selected from configured effect-modifier columns, typically baseline variables plus chronic and inflammation/sepsis proxy states when present.  The adjustment matrix \(W\) is the observed confounder set selected from the DAG helper.  The implementation can allow overlap between \(X\) and \(W\), because configured effect modifiers are not automatically removed from the adjustment set.  This is a source-supported behavior and a remaining methodological limitation.

\begin{longtable}{@{}L{0.18\textwidth}L{0.24\textwidth}L{0.24\textwidth}L{0.24\textwidth}@{}}
\caption{Implemented causal estimators and interpretation limits.}
\label{tab:causal-estimator-methods}\\
\toprule
method & implementation summary & main outputs & interpretation boundary \\
\midrule
\endfirsthead
\toprule
method & implementation summary & main outputs & interpretation boundary \\
\midrule
\endhead
Greedy Hamming matching & Binary adjustment matrix, median thresholding for non-binary covariates, median imputation, nearest available control within allowed Hamming distance, no default control reuse. & Matched pairs, pair effects, match rate, final distance, and per-treatment/global summaries. & Matched descriptive contrast; not automatically ATE or ATT. \\
LinearDML & EconML \texttt{LinearDML} with random-forest nuisance outcome and treatment models, binary treatment handling, configured \(W\) and \(X\). & Patient-level CATE, optional intervals, coefficients when aligned, model artifact, global summaries. & Model-based heterogeneity estimate under project DAG, nuisance-model, overlap, and proxy-exposure assumptions. \\
CausalForest\allowbreak DML & EconML \texttt{CausalForestDML} with random-forest nuisance models and a causal-forest final stage. & Patient-level CATE, optional intervals, feature importance, model artifact, global summaries. & Flexible heterogeneity estimate; feature importance is a model diagnostic, not mechanism proof. \\
CausalPFN & Exploratory branch using deduplicated confounder and effect-modifier features.  It writes CATE CSVs and metadata but skips EconML-specific sensitivity and permutation stages. & Patient-level CATE and normalized CATE; interval columns are unavailable and stored as missing. & Exploratory only; its primary method source, implementation version, uncertainty, and comparable diagnostic support remain unresolved. \\
\bottomrule
\end{longtable}
